% !TEX engine = xelatex
% Compile with:  tectonic -X compile docs/code.tex    (recommended)
%           or:  xelatex docs/code.tex
%
% Authored for XeLaTeX (Unicode + future Asian-locale extensibility); kept
% engine-portable so any TeX install can rebuild the PDF without changes.
% The `make code-pdf` target invokes tectonic; the Makefile is the single
% source of truth for the build command.
\documentclass[11pt,a4paper]{article}
\usepackage[margin=1in]{geometry}
\usepackage[T1]{fontenc}
\usepackage{lmodern}
\usepackage{microtype}
\usepackage{xcolor}
\definecolor{ink}{HTML}{1F1F1F}
\definecolor{slate}{HTML}{53606E}
\definecolor{rule}{HTML}{D6D6D6}
\definecolor{good}{HTML}{4F6B52}
\definecolor{bad}{HTML}{9E5B4F}
\definecolor{codebg}{HTML}{F4F4F4}
\usepackage{listings}
\lstset{
  basicstyle=\ttfamily\small\color{ink},
  backgroundcolor=\color{codebg},
  keywordstyle=\color{slate}\bfseries,
  commentstyle=\color{slate}\itshape,
  stringstyle=\color{good},
  showstringspaces=false,
  breaklines=true,
  frame=single,
  rulecolor=\color{rule},
  framesep=6pt,
  xleftmargin=8pt,
  language=Python,
}
\usepackage{tcolorbox}
\tcbuselibrary{breakable, skins}
\newtcolorbox{tip}{colback=codebg, colframe=slate, left=8pt, right=8pt, top=4pt,
  bottom=4pt, sharp corners, boxrule=0.5pt, breakable, fonttitle=\bfseries,
  title=Guide for coding large systems}
\usepackage{hyperref}
\hypersetup{colorlinks=true, linkcolor=slate, urlcolor=slate}
\usepackage{titlesec}
\titleformat*{\section}{\Large\bfseries\color{ink}}
\titleformat*{\subsection}{\large\bfseries\color{ink}}
\titleformat*{\subsubsection}{\bfseries\color{slate}}

\title{\bfseries Agent Skill Firewall \\ \large Source-Code Companion (\texttt{code.pdf})}
\author{}
\date{}

\begin{document}
\maketitle

\noindent
This document accompanies the \texttt{agent-skill-firewall} pip package. Each
section describes one source file: its purpose, the public API, the design
choices behind it, and the lessons it embodies for building larger
systems. Sidebars marked \emph{Guide for coding large systems} extract
those lessons explicitly so the document doubles as a teaching reference.

\tableofcontents
\bigskip\hrule\bigskip

\section{Package layout (one paragraph)}

\noindent\texttt{firewall/} is the public package. \texttt{cli.py} is the
\texttt{firewall} command-line entry point. \texttt{config.py} owns the
TOML configuration -- the firewall's policy lives here, not in the code,
which is what makes the system non-empirical. \texttt{runtime/} contains
the L2 verifier and the L1 containment substrate (egress proxy + canary
+ supervisor + events emitter). \texttt{gate/} contains the L0 static
scanner. \texttt{skills/} is the framework-agnostic skill discoverer.
\texttt{panel/} is the Textual TUI. \texttt{docker/} ships the minimal
compose stack. \texttt{examples/} bundles two didactic skills used by
the tests. Tests live in \texttt{tests/unit/}; a minimal experiment
harness lives in \texttt{experiment\_minimal/}.

\section{\texttt{firewall/config.py} -- the policy lives here}

\subsection{What it does}
Loads, validates, and writes the TOML configuration file at
\texttt{\textasciitilde/.config/firewall/config.toml} (XDG; macOS uses
\texttt{\textasciitilde/Library/Application Support/firewall/}). The
default config is held as a literal string so a fresh install can run
without any external state.

\subsection{Key dataclasses}
\begin{itemize}
  \item \texttt{DriverConfig} -- which LLM the runtime calls when it
    needs a tie-breaker judge. The \emph{name} of the env var is stored
    in \texttt{api\_key\_env}; the key itself is read at runtime, never
    written to disk.
  \item \texttt{PolicyConfig} -- one of three modes:
    \texttt{strict}, \texttt{balanced} (default), \texttt{observe}.
  \item \texttt{EgressConfig} -- \texttt{allow\_hosts} and
    \texttt{deny\_hosts}, both wildcard-suffix matched
    (\texttt{*.example.com}).
  \item \texttt{SecretsConfig} -- filesystem paths and env-var names
    treated as SECRET sources for taint analysis.
  \item \texttt{JudgeConfig} -- a \emph{minimal} regex list of refusal
    and compliance patterns. The user is expected to extend these in
    config without touching the source.
  \item \texttt{TaintConfig} -- the six pattern lists that drive L2's
    classifier (secret sources, network sinks, exec sinks, exfil host
    substrings, pseudo-TLDs). Every signal the runtime uses is named
    here.
  \item \texttt{Config} -- the top-level container with three
    convenience methods (\texttt{host\_allowed}, \texttt{host\_denied},
    \texttt{is\_hostile\_host}) that encapsulate the wildcard match.
\end{itemize}

\subsection{Notable functions}
\begin{lstlisting}
def load_config(path=None) -> Config:
    """Load TOML from disk, applying defaults for any missing keys."""

def init_default(path=None, *, overwrite=False) -> Path:
    """Write the bundled default TOML if no file exists."""
\end{lstlisting}

\begin{tip}
\textbf{Push policy out of code.} Patterns, thresholds, and allowlists
should read from a user-editable file. Code that \emph{decides} should
not \emph{also} carry the decision data. Here, the entire
\texttt{TaintConfig} block is data; the verifier in
\texttt{runtime/firewall.py} only knows how to read it.
\end{tip}

\section{\texttt{firewall/runtime/firewall.py} -- the L2 verifier}

\subsection{What it does}
Given an intent envelope (what the skill declared) and the evidence from
one run (what the skill did), returns an \texttt{IARVerdict}. About 250
lines, policy-free.

\subsection{Key functions}
\begin{lstlisting}
def build_intent_envelope(skill: dict) -> dict:
    """Derive {name, description, allowed_tools, allow_hosts}."""

def analyze_iar(evidence, envelope, cfg, *, honeypot_aware=False) -> IARVerdict:
    """Score one run's evidence. Returns a verdict ladder:
        BENIGN     -- in-intent only
        SUSPICIOUS -- divergence, no confirmation
        MALICIOUS  -- divergence confirmed by harmful result
    """

def recommend_enforcement(verdict) -> dict:
    """Pure data: 'none' | 'tighten_egress' | 'quarantine'."""
\end{lstlisting}

\subsection{Five invariants}
The verifier preserves five invariants regardless of user config:
\begin{enumerate}
  \item \textbf{Cross-plane taint.} A secret read on the tool-call
    plane AND an out-of-policy egress on the network plane in the same
    run = malicious, regardless of how innocent the destination looks
    alone.
  \item \textbf{Pre-scan defer.} Only \texttt{CRITICAL} static findings
    hard-block; \texttt{HIGH}/\texttt{MEDIUM} defer to the L2 layer.
  \item \textbf{Detection / containment separation.} In detector mode
    (\texttt{honeypot\_aware=False}), honeypot capture and canary hits
    are recorded under \texttt{containment} but excluded from the
    verdict. Reaching the honeypot means the detector already failed.
  \item \textbf{Unknown is not guilty.} An unfamiliar host with no
    cross-plane signal is classified \texttt{low}, not
    \texttt{suspicious}. Only taint can promote it.
  \item \textbf{IAR three-plane shape.} \texttt{intent} (envelope),
    \texttt{action} (tool calls + egress \emph{attempts}, never bodies),
    \texttt{result} (taint flow / containment). The verifier never
    reads response bodies.
\end{enumerate}

\begin{tip}
\textbf{State your invariants in writing.} If you cannot list the things
a security module must do regardless of configuration, it will drift.
Invariants belong at the top of the file and in named tests; if a
refactor breaks an invariant test, that is the refactor's problem.
\end{tip}

\section{\texttt{firewall/runtime/egress\_proxy.py} -- L1 containment}

\subsection{What it does}
A tiny HTTP/HTTPS forward proxy. Allowed hosts pass through, denied or
hostile hosts get \texttt{403}, anything else is \emph{captured}
(short-circuited with a \texttt{204} after the body is recorded as
evidence). Every decision is emitted as a JSONL event the supervisor
and the panel consume.

\subsection{Why this is L1 and not L2}
Containment captures \emph{do not} promote a verdict on their own; that
discipline lives in the verifier (invariant 3). Keeping the proxy
"dumb" -- it only knows the allow/deny lists in \texttt{cfg.egress} --
means a user can audit the containment behaviour without reading any
detection logic.

\begin{tip}
\textbf{Containment and detection belong in different files.} The egress
proxy can be replaced by an nftables rule, an Envoy sidecar, or a host
firewall; nothing in \texttt{runtime/firewall.py} cares. The events
emitter is the seam.
\end{tip}

\section{\texttt{firewall/runtime/canary.py} -- honeypot tripwire}

A 28-line HTTP server on \texttt{127.0.0.1:8088}. Any inbound request
is logged as a \texttt{CANARY-HIT}. The L2 verifier records this under
\texttt{containment} only.

\section{\texttt{firewall/runtime/supervisor.py} -- log $\rightarrow$ evidence}

\subsection{What it does}
Parses the proxy's and canary's log lines into typed \texttt{Event}
records, then aggregates a list of events into the dict shape
\texttt{analyze\_iar} expects. Keeps the regex set out of the verifier.

\section{\texttt{firewall/runtime/events.py} -- JSONL emitter}

\subsection{What it does}
A thread-safe \texttt{EventEmitter} that appends one JSON object per
line to \texttt{events.jsonl}, with size-based rotation. This is the
seam between every L1/L2 component and the panel; new event kinds can
be added without breaking readers.

\begin{tip}
\textbf{Make one wire format the contract.} The events file is the
firewall's wire format. The proxy writes it, the verifier writes it,
the panel reads it, and the experiment harness reads it. New
components fit in by writing or reading the same JSONL; no other
coupling needed.
\end{tip}

\section{\texttt{firewall/gate/static\_scanner.py} -- L0}

\subsection{What it does}
A small static scanner with three rules: hidden-instruction HTML
comments in \texttt{SKILL.md}, shell-fetch lines in code, and
hostile-host literals (from \texttt{cfg.taint.exfil\_host\_substrings}).
Returns \texttt{Finding} records with a severity; never decides
install / no-install on its own (the prescan rule lives in the
verifier, invariant 2).

\section{\texttt{firewall/skills/discover.py} -- framework-agnostic loader}

\subsection{What it does}
Walks the directories every agent framework uses for skills
(\texttt{.claude/skills/}, \texttt{.openclaw/skills/},
\texttt{.cursor/skills/}, \texttt{.opencode/skills/}, \texttt{skills/})
and parses each \texttt{SKILL.md} into a \texttt{Skill} dataclass.
Optionally falls back to PyYAML for the manifest body; the built-in
parser handles the small subset that real skills use.

\subsection{Why this matters}
The Agent Skill format converges across frameworks: a directory whose
name matches a \texttt{name} field in a single text manifest. Hooking
at the filesystem prefix means the firewall is uniform across
Anthropic SDK, Claude Code, OpenClaw, Cursor, OpenCode, etc.

\begin{tip}
\textbf{Hook at the lowest common denominator.} If multiple frameworks
agree on the file-system layout, that is your interface, not the agent
API of any one of them. The number of agents will grow; the file-system
semantics will not.
\end{tip}

\section{\texttt{firewall/panel/tui.py} -- live Textual panel}

A four-region grid: passed/blocked table, session stats, API usage,
and a tailing log. Reads \texttt{events.jsonl} only -- no firewall
enforcement of its own. Keyboard: \texttt{q} quit, \texttt{p} pause,
\texttt{c} clear table.

\section{\texttt{firewall/cli.py} -- subcommands}

The Typer-based front door. Subcommands:
\texttt{config init/show/path}, \texttt{scan}, \texttt{watch},
\texttt{start} (\texttt{--mode host|docker}), \texttt{stop},
\texttt{panel}, \texttt{skills}. No subcommand carries policy; each
delegates to its module.

\section{Tests (\texttt{tests/unit/})}

Four small files:
\begin{itemize}
  \item \texttt{test\_firewall\_invariants.py} -- one test per
    invariant from the L2 verifier (a sixth test pins the
    pseudo-TLD escalation).
  \item \texttt{test\_skill\_discovery.py} -- frontmatter parser +
    isolated-dir discovery.
  \item \texttt{test\_supervisor.py} -- regex set for the egress /
    canary log lines.
  \item \texttt{test\_static\_scanner.py} -- the canary example flags
    the hostile-host literal; the safe example is clean.
\end{itemize}

\begin{tip}
\textbf{One test per invariant, named after it.} A failing test for the
cross-plane-taint rule should literally read
\texttt{test\_cross\_plane\_taint\_is\_malicious}. Invariant tests double
as documentation.
\end{tip}

\section{Docker (\texttt{firewall/docker/})}

A minimal image (\texttt{python:3.11-slim}, non-root, read-only
root-fs in compose) that runs the egress proxy on \texttt{8080} and
the canary on \texttt{8088} via a 14-line \texttt{entrypoint.sh}. The
compose file ships a commented \texttt{openclaw} service stub so the
user can wire any agent that obeys \texttt{HTTP\_PROXY}.

\section{Experiment harness (\texttt{experiment\_minimal/})}

Two files: \texttt{run.py} loops over a skills folder and prints a CSV
of \texttt{(skill, static\_severity, runtime\_verdict, score, reason)};
\texttt{panel.html} is a 90-line browser dashboard that polls
\texttt{events.jsonl}.

\section{Agent integration commands (\texttt{firewall hook}, \texttt{firewall doctor})}

\texttt{firewall hook openclaw} prints the env vars + per-provider
config commands to route OpenClaw's HTTP egress through the proxy on
\texttt{127.0.0.1:8080}. Verified against \texttt{openclaw 2026.6.10}:
\texttt{HTTP\_PROXY}, \texttt{HTTPS\_PROXY}, and \texttt{ALL\_PROXY} are
honoured at the undici layer and catch both OpenClaw's own HTTP and
any subprocess spawned by skills (curl, pip, npm). For model-provider
calls specifically, \texttt{models.providers.<id>.request.proxy.mode}
and \texttt{.proxyUrl} pin the path. The same command supports
\texttt{claude-code}, \texttt{cursor}, and \texttt{generic} targets.

\texttt{firewall doctor} runs a six-step health check: config parses,
the configured API-key env var is set, the proxy and canary ports are
free, and the two optional tools (\texttt{docker},
\texttt{tectonic}) are reachable. Exit status is \texttt{0} on success
and \texttt{1} if any required check fails. Pair it with
\texttt{firewall config show} when triaging a fresh install.

\section*{Closing notes on system design}

Three habits this codebase encodes:
\begin{enumerate}
  \item \textbf{Push policy out of code.} The runtime decides; the
    config carries the data it decides over. Patterns rot; code that
    reads them does not.
  \item \textbf{Separate detection and containment.} They have different
    failure modes. Mixing them produces a system you cannot audit when
    the user disagrees with a verdict.
  \item \textbf{One wire format is the contract.} An append-only JSONL
    is the seam. Components plug in by reading or writing it.
\end{enumerate}

\end{document}
